\section{Evaluating conduct and manager selection}

\begin{qfmethod}[Evaluating practices, policies and conduct against the Code and Standards: a sequence for multi-standard cases, and three worked cases with reasoned conclusions]\label{meth:evaluating-conduct-cases}
\qffield{Statement} A sequence for deciding whether conduct violates the Code and Standards (Derivation, from the structure of the Standards and the framework of \hyperref[meth:ethical-decision-framework]{ethical decisions}):
\begin{enumerate}
\item \emph{Who is bound and in what role?} Members and candidates are bound in their professional activities; whether a person is an employee, contractor, supervisor or adviser decides which duties apply \citep[pp.~36, 110--111, 120]{standards-practice-handbook-2024}.
\item \emph{Which law applies, and is it stricter?} (I(A)).
\item \emph{To whom are duties owed} (beneficiaries, clients, employer, market) and do any conflict? Market integrity and clients come before the employer, and both before the member (Code; pp.~xv, 110).
\item \emph{Facts and knowledge.} Many Standards turn on what the person knew or should have known (I(C), V(B)), on intent (II(B)) or on reasonable efforts (III(D), IV(C)).
\item \emph{Standard by standard}, test the facts against each requirement, including procedures that were absent.
\item \emph{Conclusion and remedy}: which Standards are violated, and what conduct would have conformed, including dissociation and prevention.
\end{enumerate}

\qffield{Practice} Three cases (original), with conclusions.

\emph{Case 1: the fund launch.} Talia, a charterholder and head of product at Eston Funds, launches a ``capital-protected'' note whose protection depends on a bank's credit. The marketing, approved by her, describes the note as ``guaranteed'' and shows a ten-year backtest as ``performance'', without saying it is simulated. Her bonus is a share of first-year sales.
\begin{itemize}
\item I(C) and III(D): \emph{violated}: the guarantee is a credit exposure, and simulated results are presented as actual.
\item V(B)(3): \emph{violated}: the significant risk, the bank's credit, is not disclosed.
\item VI(A): the sales bonus is a conflict for the product head; at minimum it should be disclosed to the employer, and the firm should separate approval of marketing from sales incentives.
\item Conforming conduct: describe the protection as conditional on the issuer's solvency, label the backtest as hypothetical, and route approval through compliance.
\end{itemize}

\emph{Case 2: the pension consultant.} Kwame, a consultant, recommends three managers to a public pension plan. One of them pays his firm for ``data services'' and sponsored his attendance at a conference abroad. He discloses neither. The plan's trustees receive a ranking based on three-year returns only.
\begin{itemize}
\item I(B) and VI(A): \emph{violated}: payments and travel from a manager under evaluation could reasonably be expected to affect his objectivity, and were not disclosed.
\item V(A): \emph{violated}: ranking on three-year returns without operational or process review is not a reasonable basis.
\item III(A): the duty runs to the plan's beneficiaries, whose interests the undisclosed relationship puts at risk.
\end{itemize}

\emph{Case 3: the policy review.} A firm's written policy allows employees to trade any stock after a 24-hour delay following client trades, and requires annual (not quarterly) holdings reports. No violation has been found.
\begin{itemize}
\item The policy does not itself violate VI(B), which requires client priority, not specific periods.
\item But for a supervisor (IV(C)) it is weak: a 24-hour delay leaves room to trade on knowledge of pending client orders, and annual reports cannot detect it. Recommend pre-clearance, a longer blackout, duplicate confirmations and quarterly reporting (pp.~165--166).
\end{itemize}

\qffield{Conditions} The sequence is an organising device, not a CFA Institute procedure; conclusions follow the Handbook's guidance on each Standard.

\qffield{Pitfalls} Stopping at the first violation found. Treating the absence of a detected violation as proof that a policy is adequate. Ignoring who the client is when beneficiaries differ from the hiring party.

\qffield{See also} \hyperref[def:code-of-ethics-standards]{the Code and Standards}; \hyperref[meth:standards-in-manager-selection]{manager selection}.

\end{qfmethod}

\begin{qfmethod}[Applying the Code and Standards to manager selection: duties of committee members, staff and consultants hiring managers, and of managers seeking mandates, in an endowment case]\label{meth:standards-in-manager-selection}
\qffield{Statement} Manager selection involves people on both sides, and the Standards apply to each \citep[pp.~13, 132--133]{standards-practice-handbook-2024}:
\begin{itemize}
\item \emph{Hiring side} (committee members, staff, consultants): do not solicit gifts, contributions or other benefits that could influence the decision, including donations to a favoured charity or political campaign (I(B)); review external managers as diligently as securities, with consistent, objective criteria: the manager's code of ethics, compliance and controls, the quality of performance information, and adherence to strategy (V(A)). Possible criteria are the Asset Manager Code, GIPS and model requests for proposal (p.~133). Duties run to the beneficiaries (III(A)); personal or business relationships with candidates are conflicts to avoid or disclose (VI(A)).
\item \emph{Selling side} (managers): offer nothing that could impair the decision-maker's independence (I(B)); present performance fairly and completely (I(C), III(D)); disclose services, costs, process and risks (V(B)); disclose referral or placement fees (VI(C)).
\item Supervisors must ensure the process is followed (IV(C)).
\end{itemize}
``Pay-to-play'' scandals, in which managers donated to the campaigns of officials who chose managers for public pension funds, led to reporting requirements and bans on hiring contributors (p.~13).

\qffield{Practice} \emph{The Westbrook University endowment} (original). The investment committee is searching for a global small-cap manager. Facts:
\begin{enumerate}
\item The chief investment officer, Paula, a charterholder, attended a two-day conference in Lisbon paid for by Candidate~A, which also sponsors the university's business-school case competition.
\item A committee member, Hugh, is a former partner and current investor in Candidate~B.
\item The consultant scored managers on five-year returns and fees only; Candidate~C's returns are gross of fees and its composite excludes closed accounts.
\item Candidate~D offered to donate USD~50,000 to the athletics programme ``in appreciation of the partnership'' if hired.
\end{enumerate}
Evaluation:
\begin{itemize}
\item Paula: the paid trip from a candidate under evaluation \emph{violates} I(B) unless the endowment repays it; she should recuse herself from scoring A or disclose the trip to the committee, and the endowment should pay for future due-diligence travel. The sponsorship of the case competition is an institutional relationship to disclose (VI(A)).
\item Hugh: his investment in B is a conflict; he should recuse himself from the vote. If he is a member or candidate, failing to disclose it \emph{violates} VI(A).
\item The consultant's process \emph{lacks a reasonable basis} (V(A)): no operational due diligence, no check of composite construction. C's presentation, if it implies net and complete results, \emph{violates} I(C) and III(D) for any member responsible.
\item D's offer is an inducement to the hiring decision; if made by a member, it \emph{violates} I(B). The committee should disqualify D or at least refuse the offer and record it.
\item Remedy: a written selection policy with conflict declarations before each search, recusal rules, endowment-paid travel, a request for proposal that requires GIPS-compliant composites and the managers' codes and compliance manuals, and a log of contacts with candidates.
\end{itemize}

\qffield{Conditions} Applies to members and candidates; public pension plans are also subject to pay-to-play and procurement law.

\qffield{Pitfalls} Treating gifts to the institution, rather than to the individual, as harmless. Delegating selection to a consultant without reviewing the consultant's own conflicts.

\qffield{See also} \hyperref[meth:manager-due-diligence]{manager due diligence}; \hyperref[def:asset-manager-code]{Asset Manager Code}; \hyperref[meth:endowment-liquidity-allocation]{endowment liquidity}.

\end{qfmethod}
